%=============================================================================!
% 11.0  CHAPTER OPENING                                                       !
%=============================================================================!
The code of Chapter~10 advances every atom with the same step. Chapter~9
showed why that step must be short compared with the fastest vibration:
in water, an oxygen--hydrogen bond stretches with a period of about
\SI{9}{\femto\second}, restricting the step to around
\SI{0.5}{\femto\second}. Molecular rotation and rearrangement take
picoseconds, so resolving the bond vibration requires thousands of steps
to follow these slower motions.

We examine three ways to reduce that cost. Holding bonds at fixed lengths
removes their vibrations; SHAKE and RATTLE calculate the constraint forces
needed at each step. Splitting fast and slow forces allows the costly
slow forces to be evaluated less often, although resonance limits the
outer step. Moving some mass from the heavy atoms to their hydrogens
slows the fastest vibrations without changing the potential energy. We
compare these choices on a box of water and use its trajectories to
identify their limits and common implementation faults. The constraints
also change the number of degrees of freedom, which Chapter~12 needs to
relate kinetic energy to temperature.
